% !TeX root = main.tex
% English abstract and keywords; synchronized with 00_摘要.md.
\chapter*{Abstract}
\addcontentsline{toc}{chapter}{Abstract}

Mobile manipulators combine mobility and dexterity to extend robotic workspaces for indoor transport and articulated object manipulation. In dynamic environments, obstacle motion, manipulation reachability, and arm tracking lag jointly affect tracking quality and task efficiency. This thesis investigates predictive navigation and constrained manipulation through a control and evaluation pipeline incorporating navigation, whole-body coordination, and visual guidance.

The navigation system combines geometric model predictive control using \(SE(2)\) tracking errors, predictive barrier constraints, and a protection layer driven by raw laser scans. For coordinated control, arm position setpoints and their tracking dynamics are included in prediction, allowing base and arm commands to be optimized jointly. Visual measurements provide approach and grasp references.

In existing simulations, the navigation method reached the goal on all 83 routes jointly evaluated across three methods, with contact on one route. Three paired drawer experiments showed task-time reductions of 13.38--14.05\% relative to a strict stationary-base strategy, with both strategies meeting the original task criteria. A separate three-pair study found lower end-effector position errors over the full control window and fewer arm-command reversals with five-step prediction than with one-step prediction.

Visual integration achieved moving pick-and-place in specified single-bottle scenarios, while excessive grasp slip remained in a pedestrian-interaction case. One integrated door trial completed visual approach, opening, and withdrawal under ideal grasp binding and non-real-time computation. These results support feasibility under the tested simulation conditions. Complete residential retrieval, dynamic whole-body avoidance, and real-time hardware execution remain to be validated.

\thesiskeywords{Keywords: }{mobile manipulators, geometric model predictive control, control barrier functions, whole-body coordination, manipulation on the move, visual guidance}
